\subsection{正则表示}

设 \(p\) 是满足 \(\geqslant 1\) 的实数，\(\mu'\) 是 \(G\) 上的右哈尔测度。对于每个 \(g \in G\) 及每个函数 \(f \in \mathcal{L}^p(G, \mu)\)，记 \(\gamma_G^{(p)}(g)f\) 为函数 \(x \mapsto f(g^{-1}x)\)（定义在 \(G\) 上）。映射 \(g \mapsto \gamma_G^{(p)}(g)\) 是 \(G\) 在 \(\mathcal{L}^p(G)\) 中的连续线性表示。通过取商，它诱导出 \(G\) 在 \(L^p(G)\) 中的连续等距线性表示（INT, VIII, p. 135, § 2, no. 5, prop. 8），记为 \(\gamma_G^{(p)}\)。

同样，对于每个 \(g \in G\) 及每个函数 \(f \in \mathcal{L}^{p}(G, \mu')\)，记 \(\boldsymbol{\delta}_{\mathrm{G}}^{(p)}(g)f\) 为 G 上的函数 \(x \mapsto f(xg)\)。映射 \(g \mapsto \boldsymbol{\delta}_{\mathrm{G}}^{(p)}(g)\) 是 G 在 \(\mathcal{L}^{p}(G, \mu')\) 中的连续线性表示。通过取商，它诱导出 G 在 \(\mathrm{L}^{p}(\mathrm{G}, \mu')\) 中的连续等距线性表示（参见 INT, VIII, p. 136, § 2, no. 5）。

称 \(\boldsymbol{\gamma}_{\mathrm{G}}^{(p)}\) 为 G 在 \(\mathcal{L}^{p}(\mathrm{G})\) 或 \(\mathrm{L}^{p}(\mathrm{G})\) 中的左正则表示，称 \(\boldsymbol{\delta}_{\mathrm{G}}^{(p)}\) 为 G 在 \(\mathcal{L}^{p}(\mathrm{G},\mu^{\prime})\) 或 \(\mathrm{L}^{p}(\mathrm{G},\mu^{\prime})\) 中的右正则表示。

引理 3. — 设 p 是大于等于 1 的实数。G 在 L \(^{p}\) (G,μ) 中的左（分别为右）正则表示（分别在 L \(^{p}\) (G,μ') 中）是忠实的。

更确切地，令 q 为 p 的共轭指数，令 g 为 G 中满足 \(g \neq e\) 的元素。

\begin{enumerate}
\def\labelenumi{\alph{enumi})}
\item
  存在函数 \(\varphi\in\mathcal{K}(\mathrm{G})\)，它在 \(\mathrm{L}^{q}(\mathrm{G},\mu)\) 中为正且非零，并满足 \(\langle\varphi,\boldsymbol{\gamma}_{\mathrm{G}}^{(p)}(g)\varphi\rangle=0\)；
\item
  存在函数 \(\varphi\in\mathcal{K}(\mathrm{G})\)，它在 \(\mathrm{L}^{q}(\mathrm{G},\mu')\) 中为正且非零，并满足 \(\langle\varphi,\boldsymbol{\delta}_{\mathrm{G}}^{(p)}(g)\varphi\rangle=0\)。
\end{enumerate}

考虑左正则表示 \(\boldsymbol{\gamma}_{\mathrm{G}}^{(p)}\) 的情形，右正则表示的情形类似。断言 \(a\) 蕴含 \(\boldsymbol{\gamma}_{\mathrm{G}}^{(p)}\) 是忠实的，因为若 \(g \neq e\) 是 G 中的元素，且 \(\varphi\) 如 \(a\) 所述，则有 \(\varphi \neq \boldsymbol{\gamma}_{\mathrm{G}}^{(p)}(g)\varphi\)，因为 \(\langle \varphi, \varphi \rangle = \int_{\mathrm{G}} \varphi^2 > 0\)。

因此证明 \(a\)。令 \(g \neq e\) 属于 G。令 C 是 \(e\) 的一个紧对称邻域，使得 \(g \notin \mathrm{C}^2\)，并令 \(\varphi \in \mathcal{K}(\mathrm{G})\) 是积分为 1 且支撑包含于 C 的连续正函数；函数 \(\varphi\) 在 \(\mathrm{L}^p(\mathrm{G}, \mu)\) 中非零。由于 \(\mathrm{C} \cap g^{-1}\mathrm{C} = \emptyset\)，有

\[
\langle \varphi , \boldsymbol {\gamma} _ {\mathrm{G}} ^ {(p)} (g) \varphi \rangle = \int_ {\mathrm{G}} \varphi (x) \varphi (g ^ {- 1} x) d \mu (x) = 0.
\]

设 \(\varrho\) 是 G 在巴拿赫空间 E 中的连续有界表示。对于每个 \(f\in \mathrm{L}^1 (\mathrm{G})\)（分别为 \(f^{\prime}\in \mathrm{L}^{1}(\mathrm{G},\mu^{\prime}))\)）以及每个 \(g\in \mathrm{G}\)，有

\[
\varrho (g) \varrho (f \cdot \mu) = \varrho (\varepsilon_ {g} * (f \cdot \mu)) = \varrho (\boldsymbol {\gamma} _ {\mathrm{G}} ^ {(1)} (g) f \cdot \mu),\tag{1}
\]

\[
\varrho (f ^ {\prime} \cdot \mu^ {\prime}) \varrho (g) = \varrho ((f ^ {\prime} \cdot \mu^ {\prime}) * \varepsilon_ {g}) = \varrho (\boldsymbol {\delta} _ {\mathrm{G}} ^ {(1)} (g ^ {- 1}) f ^ {\prime} \cdot \mu ^ {\prime}),\tag{2}
\]

（INT, VIII, p. 144, § 3, no. 2, 公式（5））。

若 G 是单模群，置 \(\mu' = \mu\)。G 在 \(\mathcal{L}^{p}(\mathrm{G})\)（分别在 \(\mathrm{L}^{p}(\mathrm{G})\)）中的双正则表示 \(\varrho_{\mathrm{G}}^{(p)}\)，是 \(G \times G\) 在 \(\mathcal{L}^{p}(\mathrm{G})\)（分别在 \(\mathrm{L}^{p}(\mathrm{G})\)）中的表示，满足

\[
\boldsymbol {\varrho} _ {\mathrm{G}} ^ {(p)} (g _ {1}, g _ {2}) = \boldsymbol {\gamma} _ {\mathrm{G}} ^ {(p)} (g _ {1}) \boldsymbol {\delta} _ {\mathrm{G}} ^ {(p)} (g _ {2}) = \boldsymbol {\delta} _ {\mathrm{G}} ^ {(p)} (g _ {2}) \boldsymbol {\gamma} _ {\mathrm{G}} ^ {(p)} (g _ {1}).
\]

这是连续线性表示（V, p. 377 的引理 1）。它满足

\[
(\boldsymbol {\varrho} _ {\mathrm{G}} ^ {(p)} (g _ {1}, g _ {2}) f) (x) = f (g _ {1} ^ {- 1} x g _ {2})
\]

对所有 \(f \in \mathcal{L}^p(\mathrm{G})\)、所有 \((g_1, g_2) \in \mathrm{G} \times \mathrm{G}\) 及每个 \(x \in \mathrm{G}\) 成立。

注. --- G 在 \(\mathrm{L}^{p}(\mathrm{G},\mu)\) 中的双正则表示不一定忠实；其核是 G 的中心在映射 \(g\mapsto(g,g)\) 下的像（V, p. 487 的习题 4）。

当 \(p = 2\) 时，G 的左正则表示 \(\gamma_{\mathrm{G}}^{(2)}\) 在复希尔伯特空间 \(\mathrm{L}^2 (\mathrm{G},\mu)\) 中是酉的，因为它是等距的。同样，右正则表示 \(\delta_{\mathrm{G}}^{(2)}\) 在 \(\mathrm{L}^2 (\mathrm{G},\mu')\) 中也是酉的。

我们简记 \(\gamma_{\mathrm{G}} = \gamma_{\mathrm{G}}^{(2)}\) 和 \(\delta_{\mathrm{G}} = \delta_{\mathrm{G}}^{(2)}\)，并将这些表示称为 G 的左、右正则表示。

若 G 是单模群，则双正则表示 \(\varrho_{\mathrm{G}}^{(2)}\) 是 \(G \times G\) 在 \(\mathrm{L}^{2}(\mathrm{G}, \mu)\) 中的表示，且是酉的。我们简记它为 \(\varrho_{G}\)。

引理 4. --- 设 \(p\) 是满足 \(\geqslant 1\) 的实数。对于 \(f \in \mathcal{L}^1(\mathrm{G})\)，线性映射 \(\gamma_{\mathrm{G}}^{(p)}(f)\) 与 \(\mathrm{L}^p(\mathrm{G})\) 上由 \(\varphi \mapsto f *^\mu \varphi\) 定义的自同态重合。

这由 INT, VIII, p. 157, § 4, n°2 的命题 6 得出，同时要注意 INT, VIII, p. 165 的公式（14）。

注. --- 回忆一下，对于 G 上的任意函数 \(f\)，在 G 上定义函数 \(\check{f}\)，使得 \(\check{f}(g) = f(g^{-1})\) 对所有 \(g \in \mathrm{G}\) 成立（INT, VII, p. 12, § 1, n° 1, 公式（12））。

可以检验，若 \(f \in \mathcal{L}^{1}(\mathrm{G})\)，线性映射 \(u = \boldsymbol{\delta}_{\mathrm{G}}^{(p)}(f)\) 满足关系式

\[
\widetilde {u (\varphi)} = f * \check {\varphi}
\]

对所有 \(\varphi\in\mathrm{L}^{p}(\mathrm{G},\mu^{\prime})\) 成立。
% semantic-fragments: legacy-input-seam
\relax{}
